\chapter{Discussion and Conclusion}
\label{ch:discussion}

This chapter synthesises findings across the six sub-experiments, draws out implications for platform design and competition policy, and closes with the main limitations and a short list of directions for future work. I treat every simulation output as a black-box observation. The discussion therefore compares the observed patterns to theoretical predictions and prior empirical work but does not assign mechanistic explanations to the simulation findings.

\section{Synthesis Across Experiments}
\label{sec:synthesis}

\Cref{tab:cross_exp_summary} gathers the top three revenue drivers, the mean revenue by auction format, and the format gap for each of the six sub-experiments. Two cross-cutting patterns stand out.

\begin{table}[H]
\centering
\caption[Cross-experiment summary of revenue drivers and format effects]{\textbf{Cross-experiment summary of revenue drivers and format effects.} For each experiment, the table reports the top three factors by absolute $t$-statistic on platform revenue, the mean revenue under first-price and second-price formats, and the percentage gap between them defined as $(\mathrm{FPA} - \mathrm{SPA}) / \mathrm{SPA}$.}
\label{tab:cross_exp_summary}
\resizebox{\textwidth}{!}{%
\begin{tabular}{l l l l l c c c}
\toprule
& & \multicolumn{3}{c}{\textbf{Top revenue drivers} ($|t|$)} & \multicolumn{3}{c}{\textbf{Revenue by format}} \\
\cmidrule(lr){3-5} \cmidrule(lr){6-8}
\textbf{Exp} & \textbf{Algorithm} & \textbf{Rank 1} & \textbf{Rank 2} & \textbf{Rank 3} & \textbf{FPA} & \textbf{SPA} & \textbf{Gap (\%)} \\
\midrule
1a & Q-learning constant & \ExpOneARevFactorOne{} (\ExpOneARevFactorTOne) & \ExpOneARevFactorTwo{} (\ExpOneARevFactorTTwo) & \ExpOneARevFactorThree{} (\ExpOneARevFactorTThree) & \ExpOneARevMeanFPA & \ExpOneARevMeanSPA & $\ExpOneARevGapPct$ \\
1b & Q-learning affil.\ & \ExpOneBRevFactorOne{} (\ExpOneBRevFactorTOne) & \ExpOneBRevFactorTwo{} (\ExpOneBRevFactorTTwo) & \ExpOneBRevFactorThree{} (\ExpOneBRevFactorTThree) & \ExpOneBRevMeanFPA & \ExpOneBRevMeanSPA & $\ExpOneBRevGapPct$ \\
2a & LinUCB & \ExpTwoARevFactorOne{} (\ExpTwoARevFactorTOne) & \ExpTwoARevFactorTwo{} (\ExpTwoARevFactorTTwo) & \ExpTwoARevFactorThree{} (\ExpTwoARevFactorTThree) & \ExpTwoARevMeanFPA & \ExpTwoARevMeanSPA & $\ExpTwoARevGapPct$ \\
2b & Thompson & \ExpTwoBRevFactorOne{} (\ExpTwoBRevFactorTOne) & \ExpTwoBRevFactorTwo{} (\ExpTwoBRevFactorTTwo) & \ExpTwoBRevFactorThree{} (\ExpTwoBRevFactorTThree) & \ExpTwoBRevMeanFPA & \ExpTwoBRevMeanSPA & $\ExpTwoBRevGapPct$ \\
3a & Dual pacing & \ExpThreeARevFactorOne{} (\ExpThreeARevFactorTOne) & \ExpThreeARevFactorTwo{} (\ExpThreeARevFactorTTwo) & \ExpThreeARevFactorThree{} (\ExpThreeARevFactorTThree) & \ExpThreeARevMeanFPA & \ExpThreeARevMeanSPA & $\ExpThreeARevGapPct$ \\
3b & PI pacing & \ExpThreeBRevFactorOne{} (\ExpThreeBRevFactorTOne) & \ExpThreeBRevFactorTwo{} (\ExpThreeBRevFactorTTwo) & \ExpThreeBRevFactorThree{} (\ExpThreeBRevFactorTThree) & \ExpThreeBRevMeanFPA & \ExpThreeBRevMeanSPA & $\ExpThreeBRevGapPct$ \\
\bottomrule
\end{tabular}}
\end{table}

The first pattern is that structural market parameters rank above algorithmic hyperparameters in every experiment. In the four unconstrained learning experiments, the number of bidders is the top factor for revenue in every case. In the two pacing experiments, the budget multiplier replaces it as the top factor, and the number of bidders drops to third or fourth place. Within each family, learning rates, discount factors, regularisation, memory decay, exploration intensity, and controller aggressiveness all fail to produce significant main effects on revenue at the magnitudes reached by competitive pressure or budget tightness. This empirical ranking points in the same direction as the classical result of \citet{BulowKlemperer1996}, who argue that adding one more bidder matters more than mechanism optimisation in static auctions. I stop short of claiming the two findings are equivalent, since their setting is a one-shot auction with rational bidders and mine is a repeated auction with algorithms, but the directional agreement is worth noting.

The second pattern is that the sign of the auction-format effect on revenue differs between the learning family and the pacing family. Under Q-learning with constant valuations, the two formats generate nearly identical revenue. Under Q-learning with affiliated valuations and under both contextual-bandit experiments, second-price auctions generate higher revenue, with format gaps of roughly 11, 19, and 9 percent respectively. Under both pacing experiments, first-price auctions generate higher revenue, with format gaps of roughly 9 percent under dual pacing and a larger positive gap under PI pacing. The between-family sign difference is not a small margin, and it holds under every robustness check reported in \cref{ch:robustness}. The format ranking that \citet{BanchioSkrzypacz2022} report from Q-learning experiments, in which first-price auctions appear more vulnerable to bid suppression than second-price, does not carry over to the pacing family in my simulations, where the opposite direction holds.

The two patterns combine into a single message for auction designers. A platform asking whether it should use first-price or second-price formats cannot read a universal recommendation off the economic-theory shelf, nor off the reinforcement-learning literature, nor off the autobidding literature. The ordering depends on the bidding technology the agents actually use. If a platform's advertisers deploy budget-constrained pacing systems similar to the ones I study, my simulation output would recommend first-price formats; if they run unconstrained learning algorithms with affiliated signals similar to mine, it would recommend second-price. Platforms running both kinds of agents face a trade-off that my experiments cannot resolve on their own, since the thesis studies each algorithm family in isolation and does not simulate mixed populations.

\section{Policy Implications}
\label{sec:policy}

Two implications follow for competition policy, subject to the caveat that both are drawn from simulation output and depend on the external validity of the laboratory. The first is that the factor with the largest measured effect on seller revenue, in every experiment of this thesis, is structural. The number-of-bidders factor ranks first in every learning experiment, and the budget multiplier ranks first in every pacing experiment. Algorithmic tuning factors never match their magnitudes. If that ranking carries over to real markets, interventions that raise the number of active bidders or that loosen budget constraints would do more for revenue than any of the algorithmic or format-based interventions I vary in the factorial designs. The corresponding policy ideas, lowering barriers to entry, preventing excessive buyer concentration, and avoiding restrictive spend caps, are ordinary structural remedies and my results give them the same directional support they already have in classical auction theory \citep{BulowKlemperer1996}.

The second implication is that format-based regulation faces a heterogeneity problem. A rule that mandates one auction format as safer than the other requires that the format ranking be stable across the bidding technologies regulators will actually encounter. My findings in \crefrange{ch:qlearning}{ch:pacing} show that, within the settings I study, it is not. A rule favouring second-price auctions on the basis of the Q-learning literature would misfire in markets where pacing bidders predominate, and a rule favouring first-price auctions on the basis of the pacing literature would misfire in markets dominated by unconstrained learners. Regulators who wish to address algorithmic bid suppression would need to condition their recommendations on the bidding technology they audit, which is harder than issuing a blanket rule. Outcome-based monitoring, in which the regulator looks at realised revenue shortfalls rather than the choice of format, sits more comfortably with this heterogeneity. \citet{Hartline2024} propose regret auditing in the same spirit, where algorithmic bidders are evaluated against the performance they could have achieved with the best fixed action rather than against claims about intent.

A related implication concerns the evidentiary threshold for antitrust liability. Both United States antitrust law and European competition law require evidence of agreement or concerted practice to establish liability for horizontal coordination, and no jurisdiction has yet sanctioned purely autonomous tacit collusion by independent learning algorithms \citep{EzrachiStucke2017, Mehra2016}. The factorial laboratory supplies one ingredient that such sanctions would need, which is a transparent and reproducible ranking of the design factors associated with revenue shortfalls, along with direction, magnitude, and statistical significance. A regulator could in principle require platforms to run standardized factorial experiments on their own bidding systems and report the ranked effects, or could construct a parallel laboratory independently. The ranked-effects tables and Sobol decompositions in this thesis illustrate the kind of output that such a compliance exercise would deliver. Whether any jurisdiction would in fact adopt this approach is a separate question that my experiments do not address.

\section{Limitations}
\label{sec:limitations}

The experiments are stylised. The bidder populations are small, symmetric, and drawn from a handful of value models. The algorithms are tabular Q-learning, linear contextual bandits, and two pacing variants; they do not include deep reinforcement learning, transformer-based bidders, or large-language-model agents. There is no explicit communication channel between bidders, no combinatorial allocation, no auction-revenue feedback into bidder populations, and no entry or exit dynamics. Episode lengths are fixed and burn-in periods are standardised. The log-normal valuation model in the pacing experiments is specifically chosen to mirror the stylised setting of the autobidding literature and is not calibrated to any real exchange. External validity to a specific market requires empirical validation with proprietary auction data, which I do not have.

The black-box framing is both a strength and a limitation. It disciplines the interpretation of results, since I report only what the input-output map shows, but it leaves unanswered the mechanistic question of why the observed patterns arise. I cannot tell whether the auction-format reversal under pacing is driven by the dual-update dynamics, the budget cap, the pacing feedback loop, or some interaction among these features. Mechanistic investigations using either analytical pacing theory or targeted ablations would complement the black-box diagnosis.

Finally, the factorial design spans a fixed set of factor levels. The rankings are valid within the ranges tested. Extreme misconfigurations, such as near-zero learning rates, implausibly large budgets, or a single dominant bidder with monopsony power, could change the ordering, but those cases are unlikely to represent any real market of economic interest.

\section{Future Directions}
\label{sec:future}

Three extensions seem particularly worthwhile. The first is to apply the same factorial framework to deep reinforcement-learning agents (neural networks trained to bid by trial and error) and to large-language-model bidders (systems that receive prompts and return bid decisions), since these are the technologies most likely to be deployed in newly developed auction markets. A replication of the format-reversal result with these algorithm families would considerably strengthen its policy relevance. The second is to move from screening designs, which identify which factors matter, to response-surface designs, which trace out the functional relationship between the factors that matter and the outcome. Knowing the shape of the relationship between the number of bidders, the budget multiplier, and revenue would allow a platform to set quantitative targets rather than simply prefer one direction to another. The third is to seek empirical validation on real auction data. Several major advertising exchanges run internal shadow experiments and factor tests, and cooperation with such a platform would let the black-box framework confront real bidding populations and check whether the laboratory rankings replicate in production.

The computational laboratory developed in this thesis offers a standardised pipeline for auction platforms and antitrust regulators who want to diagnose algorithmic behaviour without committing themselves to any particular theory of how the algorithms generate their outputs. The pipeline is portable, the analysis is mechanical, and the outputs are transparent. I hope it proves useful beyond the specific experiments reported here.

\section{A Short Plain-Language Summary}
\label{sec:plain_summary}

If I had to summarise the thesis in three sentences for a reader who wanted the economic point without the statistical machinery, they would be the following. When bidders are algorithms rather than humans, the auctioneer's revenue is shaped far more by the market's structural features (how many bidders compete and how tight their budgets are) than by the choices the algorithms' designers make about learning rates, exploration schedules, or update rules. The choice between first-price and second-price auctions does not have a universal ranking. Second-price formats raise more revenue when bidders are unconstrained learners, and first-price formats raise more revenue when bidders are budget-constrained pacers. A policy recommendation that works well for one class of algorithms can therefore be the wrong recommendation for another, and regulators, platforms, and auction designers should condition their choices on the bidding technology that actually populates the market.
